\input{preamble}

% Griffiths and Schroeter p. 388

\section*{Born approximation 2}

We will now derive the Green's function that was used in the previous section.

\bigskip
Let $G(\mathbf x)$ be a Green's function such that
\begin{equation*}
\left(\nabla^2+k^2\right)G(\mathbf x)=\delta(\mathbf x)
\tag{1}
\end{equation*}

Let $g(\mathbf y)$ be the Fourier transform of $G(\mathbf x)$ such that
\begin{equation*}
G(\mathbf x)
=\frac{1}{(2\pi)^{3/2}}
\int\exp(i\mathbf x\cdot\mathbf y)
g(\mathbf y)\,d\mathbf y
\tag{2}
\end{equation*}

By equations (1) and (2)
\begin{equation*}
\left(\nabla^2+k^2\right)
\left[\frac{1}{(2\pi)^{3/2}}\int\exp(i\mathbf x\cdot\mathbf y)g(\mathbf y)\,d\mathbf y\right]
=\delta(\mathbf x)
\end{equation*}

By linearity
\begin{equation*}
\frac{1}{(2\pi)^{3/2}}
\int\left(\nabla^2+k^2\right)\exp(i\mathbf x\cdot\mathbf y)g(\mathbf y)\,d\mathbf y
=\delta(\mathbf x)
\end{equation*}

Noting that $\nabla^2$ applies $\mathbf x$ and not $\mathbf y$ we have
\begin{equation*}
\nabla^2\exp(i\mathbf x\cdot\mathbf y)g(\mathbf y)
=-|\mathbf y|^2\exp(i\mathbf x\cdot\mathbf y)g(\mathbf y)
\end{equation*}

Hence
\begin{equation*}
\frac{1}{(2\pi)^{3/2}}
\int\left(-|\mathbf y|^2+k^2\right)\exp(i\mathbf x\cdot\mathbf y)g(\mathbf y)\,d\mathbf y
=\delta(\mathbf x)
\end{equation*}

By the identity
\begin{equation*}
\delta(\mathbf x)=\frac{1}{(2\pi)^3}\int\exp(i\mathbf x\cdot\mathbf y)\,d\mathbf y
\end{equation*}

we have
\begin{equation*}
\frac{1}{(2\pi)^{3/2}}
\int\left(-|\mathbf y|^2+k^2\right)\exp(i\mathbf x\cdot\mathbf y)g(\mathbf y)\,d\mathbf y
=\frac{1}{(2\pi)^3}\int\exp(i\mathbf x\cdot\mathbf y)\,d\mathbf y
\end{equation*}

Hence
\begin{equation*}
g(\mathbf y)=\frac{1}{(2\pi)^{3/2}\left(k^2-|\mathbf y|^2\right)}
\end{equation*}

Substitute for $g(\mathbf y)$ in (2) to obtain
\begin{equation*}
G(\mathbf x)=\frac{1}{(2\pi)^3}\int\frac{\exp(i\mathbf x\cdot\mathbf y)}{k^2-|\mathbf y|^2}\,d\mathbf y
\end{equation*}

Change to polar coordinates where $x=|\mathbf x|$ and $r=|\mathbf y|$.
%and $\theta$ and $\phi$ are the angular distance from $\mathbf x$ to $\mathbf y$.
\begin{equation*}
G(\mathbf x)=\frac{1}{(2\pi)^3}
\int_0^\infty\int_0^\pi\int_0^{2\pi}
\frac{\exp(ixr\cos\theta)}{k^2-r^2}\,r^2\sin\theta\,dr\,d\theta\,d\phi
\end{equation*}

Integrate over $\phi$ (multiply by $2\pi$).
\begin{equation*}
G(\mathbf x)=\frac{1}{(2\pi)^2}
\int_0^\infty\int_0^\pi
\frac{\exp(ixr\cos\theta)}{k^2-r^2}\,r^2\sin\theta\,dr\,d\theta
\end{equation*}

For the integral over $\theta$ use change of variable
$u=\cos\theta$ and $du=-\sin\theta\,d\theta$ to obtain
\begin{equation*}
G(\mathbf x)=\frac{1}{(2\pi)^2}
\int_0^\infty\int_1^{-1}
-\frac{\exp(ixru)}{k^2-r^2}\,r^2\,dr\,du
\end{equation*}

By the integral
\begin{equation*}
\int_1^{-1}-\exp(ixru)\,du=\frac{2\sin(xr)}{xr}
\tag{3}
\end{equation*}

we have
\begin{equation*}
G(\mathbf x)=\frac{1}{2\pi^2x}
\int_0^\infty\frac{\sin(xr)}{k^2-r^2}\,r\,dr
\end{equation*}

Noting that $r\sin(xr)$ is an even function of $r$ we can extend the lower limit as
\begin{equation*}
G(\mathbf x)=\frac{1}{4\pi^2x}
\int_{-\infty}^\infty\frac{\sin(xr)}{k^2-r^2}\,r\,dr
\end{equation*}

Negate the denominator.
\begin{equation*}
G(\mathbf x)=\frac{1}{4\pi^2x}
\int_{-\infty}^\infty-\frac{\sin(xr)}{r^2-k^2}\,r\,dr
\end{equation*}

Change the sine function to exponential form and factor the denominator.
\begin{equation*}
G(\mathbf x)
=\frac{i}{8\pi^2x}
\left[
\int_{-\infty}^\infty\frac{\exp(ixr)}{(r+k)(r-k)}\,r\,dr
-\int_{-\infty}^\infty\frac{\exp(-ixr)}{(r+k)(r-k)}\,r\,dr
\right]
\tag{4}
\end{equation*}

By Cauchy's integral formula we have
\begin{equation*}
\int_{-\infty}^\infty\frac{r\exp(ixy)}{r+k}\frac{1}{r-k}\,dr
=i\pi\exp(ikx)
\end{equation*}

and
\begin{equation*}
\int_{-\infty}^\infty\frac{r\exp(-ixy)}{r-k}\frac{1}{r+k}\,dr
=-i\pi\exp(ikx)
\end{equation*}

Hence
\begin{equation*}
G(\mathbf x)
=-\frac{\exp(ikx)}{4\pi x}
\tag{5}
\end{equation*}

where
\begin{equation*}
x=|\mathbf x|
\end{equation*}

\href{https://georgeweigt.github.io/examples/born-approximation-2-demo.html}{Eigenmath script}

\end{document}
